Fix \(P\in\mathfrak P_{\mathcal D}\) and suppose that \(\mathcal D=\{d\}\), that \(n_E,n_O\) are positive, and that
\[
 B=c_{E,d}n_E+c_{O,d}n_O.
\]
Since the unit costs and counts are strictly positive, \(B>0\).  Put
\[
 \alpha_0=\frac{c_{E,d}n_E}{B}.
\]
Both cost summands in \(B\) are positive, so \(0<\alpha_0<1\), and
\[
 1-\alpha_0
 =\frac{B-c_{E,d}n_E}{B}
 =\frac{c_{O,d}n_O}{B}.                                      \tag{1}
\]
Because \(P\in\mathfrak P_{\mathcal D}\), the sole admissible design is regular.  The regular branch of the design-value profile, equivalently the profile used in the budget-profile theorem, therefore gives
\[
 \mathcal V_d(\alpha_0;P)
 =\frac{c_{E,d}V_{E,d}}{\alpha_0}
  +\frac{c_{O,d}V_{O,d}}{1-\alpha_0}.                         \tag{2}
\]
Substituting the definition of \(\alpha_0\) and (1) into (2) yields
\[
\begin{aligned}
 \frac{1}{B}\mathcal V_d(\alpha_0;P)
 &=\frac{1}{B}\left\{
   \frac{c_{E,d}V_{E,d}}{c_{E,d}n_E/B}
   +\frac{c_{O,d}V_{O,d}}{c_{O,d}n_O/B}
   \right\}\\
 &=\frac{V_{E,d}}{n_E}+\frac{V_{O,d}}{n_O}.                  \tag{3}
\end{aligned}
\]
This is the asserted identity.  The nondegeneracy assumption makes both variance components in (3) finite and strictly positive.

By the designwise-gradient lemma, \(\mu_{d,1}-\mu_{d,0}\) is an identified observed-data functional whose source-gradient components are \(\phi_{E,d}\) and \(\phi_{O,d}\).  Each component has conditional mean zero in its own source, and
\[
 \operatorname{Var}_P(\phi_{E,d}\mid G=E)=V_{E,d},
 \qquad
 \operatorname{Var}_P(\phi_{O,d}\mid G=O)=V_{O,d}.           \tag{4}
\]
For independent source samples of sizes \(n_E\) and \(n_O\), the gradient of the two-sample functional is the sum of the two empirical-average gradients.  Independence, conditional centering, and (4) imply that its variance is
\[
 \frac{V_{E,d}}{n_E}+\frac{V_{O,d}}{n_O},                    \tag{5}
\]
which equals (3).  This proves the fixed-system pathwise-gradient variance assertion directly for the exact source scores defined in this note; it does not identify them with any differently normalized score printed in the anchor.

For the final conditional assertion, additionally suppose that \(\mu_{d,1}-\mu_{d,0}=\tau\) and that \((\phi_{E,d},\phi_{O,d})\) is the canonical gradient for \(\tau\) in the stated fixed-system model.  By the defining Hilbert-space characterization of the canonical gradient, the local semiparametric efficiency bound is the squared norm of that gradient in the product two-source experiment.  Under independent sampling with fixed source counts, this squared norm is exactly (5), and hence, by (3), it is
\[
 \frac{1}{B}\mathcal V_d\!\left(\frac{c_{E,d}n_E}{B};P\right).
\]
The additional bridge-existence, bridge-uniqueness, operator-bijectivity, moment, and interior sampling-ratio conditions delimit the regime in which the assumed causal equality and canonical-gradient statement apply.  If either the equality with \(\tau\) or canonicality is absent, the designwise-gradient lemma yields only the bridge-functional variance identity (5), so no causal efficiency conclusion is drawn.